\section{Categories and functors}


\begin{exercise}{1 (Micro)}
For $P = M \times N$ (the usual cartesian product of sets): 
Verify in detail that the map does factor uniquely. 
That is, prove that there is some $f$ such that $\mu' = \mu \circ f$ and $\nu' = \nu \circ f$, and that such an $f$ is unique.
\end{exercise}
\begin{proof}
Let $f(x) = (\mu'(x), \nu'(x))$, we then have $(\mu \circ f)(x) = \mu((\mu'(x), \nu'(x)) = \mu'(x)$ so that $\mu' = \mu \circ f$, the proof that $\nu' = \nu \circ f$ holds analogously.

To see $f$ is unique, suppose there is a $g$ such that $\mu' = \mu \circ g$ and $\nu' = \nu \circ g$, this implies that both the first and second coordinate of $g(x)$ and $f(x)$ are the same (since they map to a specific element in $M$ and $N$.
Since this holds true for all $x \in P'$, then $f = g$ proving the uniqueness of $f$.
\end{proof} 

\begin{exercise}{2 (Micro)}
Take a theory you know well, such as groups, rings, topological spaces, vector spaces. 
You have no doubt also define products here. 
Verify that the same property talked about in this page also holds in this context. 
For example, in vectors spaces, the product of two vector spaces $V \times W$ has two linear maps $\mu: V \times W \to V$ and $\nu: V \times W \to W$ such that for any linear maps $\mu': P \to V$ and $\nu': P \to W$ , there is a unique linear map $f: P \to V \times W$ such that $\mu' = \mu \circ f$ and $\nu' = \nu \circ f$.
Furthermore, show that the argument of 1.0.1 indeed shows that products are unique up to isomorphism.
\end{exercise}
\begin{proof}
We will take the category of groups with homomorphisms as its maps.
Let $(P = (G \times G'), \mu, \nu)$ and $(P', \mu', \nu')$ be given with $P'$ an arbitrary group.
Consider the function $f: P' \to P$ given by $f(x) = (\mu'(x), \nu'(x))$.
We have to prove $f$ is a homomorphism.
To see this is the case, let $x, y \in P'$, we have 
\begin{align*}
    f(x y) 
    =& (\mu'(x y), \nu'(x y))\\
    =& (\mu'(x) \mu'(y), \nu'(x) \nu'(y))\\
    =& (\mu'(x), \nu'(x)) (\mu'(y), \nu'(y))
    = f(x) f(y),
\end{align*}
where the second equality follows because $\mu'$ and $\nu'$ are homomorphisms.

To show products are unique up to isomorphism, let $(P, \mu, \nu)$ and $(P', \mu', \nu')$ be products.
Then there exists a unique homomorphism $f: P' \to P$ such that $\mu' = \mu \circ f$ and $\nu' = \nu \circ f$ and a unique homomorphism $g: P \to P'$ such that $\mu = \mu' \circ g$ and $\nu = \nu' \circ g$.
Thus, $f \circ g = id_P$ and $g \circ f = id_{P'}$;
meaning that $f$ and $g$ are invertible homomorphisms: an isomorphism between $P$ and $P'$.
\end{proof} 

\begin{exercise}{3 (Micro)}
List as many (interesting) categories (ie objects + set of morphisms) that you can, and find what the isomorphisms are in each 
\end{exercise}
\begin{proof}
Vector spaces: Linear maps; Groups: Group homomorphisms; Rings: Ring homomorphisms; Topological spaces: Continuous functions; Sets: Set functions...
\end{proof} 

\begin{exercise}{4 (Micro)}
Show that the identity morphism is unique.
\end{exercise}
\begin{proof}
Let $A, B$ be elements of a category and $id_A, id_A'$ be identities of $A$.
We have $id_A = id_A \circ id_A' = id_A'$ so that the identity is unique.
\end{proof} 

\begin{exercise}{5 (Micro)}
Show that the inverse of a morphism is unique.
\end{exercise}
\begin{proof}
Let $f: A \to B$ be a morphism and $h^{-1}: B \to A$, $g^{-1}: B \to A$ be inverses of $f$.
We have $h = h \circ id_A = h \circ f \circ g = id_A \circ g = g$, as required.
\end{proof} 

\begin{exercise}{1.1.A}
A category in which each morphism is an isomorphism is called a grupoid.
(This notion is not important in what we will discuss. 
The point of this exercise is to give you some practice with categories by relating them to an object you know well.)
\begin{enumerate}
    \item A perverse definition of a group is: a groupoid with one object.
    Make sense of this.
    (Similarly, in case you care, a perverse definition of a monoid is: a category with one object).
    \item Describe a groupoid that is not a group.
\end{enumerate}
\end{exercise}
\begin{proof}
\begin{enumerate}
    \item The point is that most of the heavy lifting of the definition of a category is about the morphisms between the objects in the category.
    Thus, the objects of the category is a single object.
    Generically, the set (class) of objects is: $\set{\ast}$, moreover, the morphisms are $Hom(\ast, \ast) = G$, the group we are defining as the category.
    For any $g \in G$, we have that $g: \ast \to \ast$.
    Note that the axioms of a group, in particular associativity and the identity element guarantee that $G$ is indeed a category.
    \item fill
\end{enumerate}
\end{proof} 

\begin{exercise}{1.1.B}
If $A$ is an object in a category $\cC$, show that the invertible elements of $Hom(A, A)$ form a group (called the automorphism group of $A$, denoted by $Aut(A)$).
What are the automorphism groups of the objects in Examples 1.1.2 and 1.1.3?
Show that two isomorphic objects have isomorphic automorphism groups.
\end{exercise}
\begin{proof}
The fact that $Aut(A)$ is a group follows directly from the axioms of categories in addition to the fact that we are taking the invertible elements of $Hom(A, A)$, which gives us the axioms of groups that is missing in the axioms of a category.

In the category of sets, the group of automorphisms are simply the set of invertible functions of a set, the set of permutations.
In the category of vector spaces, the group of automorphisms is the set of invertible linear operators.

To see that two isomorphic objects have isomorphic automorphisms groups, let $A$ and $B$ be isomorphic, so that there exists a bijection $f: A \to B$ (with inverse $f^{-1}$) between them.
I claim that the function $F: Aut(A) \to Aut(B)$ defined by $g \mapsto f \circ g \circ f^{-1}$ is a bijection.
Suppose 
\begin{align*}
    f \circ g \circ f^{-1}
    &= F(g)\\
    &= F(g') = f \circ g' \circ f^{-1},
\end{align*}
thus 
\begin{align*}
    g 
    =& f^{-1} \circ f \circ g \circ f^{-1} \circ f\\
    =& f^{-1} \circ f \circ g' \circ f^{-1} \circ f = g',
\end{align*}
so that $F$ is injective.
Let $h \in Aut(B)$, we have that $f^{-1} \circ h \circ f \in Aut(A)$, so that $F(f^{-1} \circ h \circ f) = h$, and $F$ is surjective.
Since $F$ is both injective and surjective, it is a bijection and thus $Aut(A)$ and $Aut(B)$ are isomorphic.
\end{proof} 

\begin{exercise}{6 (Micro)}
Show that the antisymmetry axiom of a partially ordered set is not necessary for the category entailed by the order to be a valid category.
\end{exercise}
\begin{proof}
Let $(S, \geq)$ be a partially ordered set without the antisymmetry assumption.
Let $S$ be the category defined by it with $x, y \in S$ and $Hom(x, y) = f$ if $x \geq y$ and $Hom(x, y) = \emptyset$ otherwise.
For $w, x, y, z \in S$, and $f \in Hom(w, x)$, $g \in Hom(x, y)$, and $h \in Hom(y, z)$, we have $w \geq x \geq y \geq z$, so that $h \circ g \in Hom(x, z)$ and $g \circ f \in Hom(w, y)$.
Thus, $(h \circ g) \circ f = h \circ (g \circ f)$, so that associativity holds.
Naturally for every element there is an identity $id_x$, since $x \geq x$ and $x \leq x$, and the identity composes in the natural way.
Note that we haven't used the antisymmetry property to prove the axioms of a category, so that this property is not necessary for the category to be valid.
\end{proof} 

\begin{exercise}{7 (Micro)}
What is the product in the category entailed by a partially ordered set?
\end{exercise}
\begin{proof}
Let $(p, \mu, \nu)$ be a product in $\C$, so that $p \geq x$ and $p \geq y$. 
Then for any other product $(p', \mu', \nu')$ so that $p' \geq x$ and $p' \geq y$ and there exists a unique morphism from $p'$ to $p$, which in this category translates to $p' \geq p$.
Thus, $p = \sup(x, y)$, the lower upper bound of $x$, and $y$, given that any other $p'$ which is an upper bound of $x$ and $y$, is greater than $p$.
\end{proof} 

\begin{exercise}{8 (Micro)}
Verify example 1.1.14 in detail.
\end{exercise}
\begin{proof}
Let $A$ be an object in a category $\cC$.
Then there is a functor $h^A: \cC \to Sets$ sending $B \in \cC$ to $Hom(A, B)$ and sending $f: B_1 \to B_2$ to $Hom(A, B_1) \to Hom(A, B_2)$, described by $[g: A \to B_1] \mapsto [f \circ g: A \to B_1 \to B_2]$.

Let $B \in \cC$ and $id_B$ be the identity of $B$. 
We have $F(id_B)$ is defined by 
\begin{align*}
    [g: A \to B] \mapsto [id_B \circ g: A \to B \to B] = [g: A \to B],
\end{align*}
the identity on $Hom(A, B)$.

Now, let $f \in Hom(B_1, B_2)$ and $h \in Hom(B_2, B_3)$ we define $F(h \circ f)$ by 
\begin{align*}
    [g: A \to B] \mapsto [(h \circ f) \circ g: A \to B_1 \to B_2 \to B_3].
\end{align*}
We know that $F(f)$ is mapped to 
\begin{align*}
    [g: A \to B_1] \mapsto [f \circ g: A \to B_1 \to B_2],
\end{align*}
and 
Furthermore, we have seen that $F(h)$ is mapped to 
\begin{align*}
    [g: A \to B_2] \mapsto [h \circ g: A \to B_2 \to B_3].
\end{align*}
So that for any map $g: A \to B_1$, $F(h) \circ F(f)$ maps to $F(h) \circ (f \circ g) = h \circ f \circ g$, which is precisely $F(h \circ f)$, so that the functor is associative.
\end{proof}


\begin{exercise}{9 (Micro)}
A group is a category with one object, take two such cats, describe the functors between these cats
\end{exercise}
\begin{proof}
The functors between two categories defined by groups are the homomorphisms between the groups that define them.
In slightly more detail, let $\cG$ and $\cG'$ be categories defined by the groups $G$ and $G'$.
As we have seen before, $Obj(\cG) = \set{\ast}$ and $Hom(\ast, \ast) = G$ and likewise for $\cG'$.
Then letting $F$ be any homomorphism between $G$ and $G'$, we have that for $g, h \in G$, $F(g, h) = F(g) F(h)$ and $F(id_\ast) = F(e) = id_{\ast'}$ (both results by the properties of homomorphisms).
\end{proof} 

\begin{exercise}{10 (Micro)}
Take two posets and consider them as cats, describe the functor between these cats
\end{exercise}
\begin{proof}
Let $P$ and $P'$ be posets, and $\cP$ an $\cP'$ the categories entailed by them.
For $id_x$, we have that $F(id_x) = id_{F(x)}$, since $F(x) = F(x)$ in $\cP'$.
Furthermore, if $x, y \in \cP$ and $x \geq y$ (so that $Hom(x, y)$ is not empty), then $F(x) \geq F(y)$ (so that $Hom(F(x), F(y))$ is not empty).
This is the set of monotonic increasing functions.
\end{proof} 

1) 
2) 

\begin{exercise}{11 (Micro)}
Verify 1.1.20 in detail
\end{exercise}
\begin{proof}
The analysis is similar to that of 1.1.4:

Let $A$ be an object in a category $\cC$.
Then there is a contravariant functor $h_A: \cC \to Sets$ sending $B \in \cC$ to $Hom(B, A)$ and sending $f: B_1 \to B_2$ to $Hom(B_2, A) \to Hom(B_1, A)$, described by $[g: B_2 \to A] \mapsto [g \circ f: B_1 \to B_2 \to A]$.

Let $B \in \cC$ and $id_B$ be the identity of $B$. 
We have $F(id_B)$ is defined by 
\begin{align*}
    [g: B \to A] \mapsto [g \circ id_B: B \to A] = [g: B \to A],
\end{align*}
the identity on $Hom(B, A)$.

Now, let $f \in Hom(B_1, B_2)$ and $h \in Hom(B_2, B_3)$ we define $F(h \circ f)$ by 
\begin{align*}
    [g: B_3 \to A] \mapsto [g \circ (h \circ f) : B_1 \to B_2 \to B_3 \to A].
\end{align*}
We know that $F(f)$ is mapped to 
\begin{align*}
    [g: B_1 \to A] \mapsto [g \circ f: B_2 \to B_1 \to A],
\end{align*}
and 
Furthermore, we have seen that $F(h)$ is mapped to 
\begin{align*}
    [g: B_2 \to A] \mapsto [g \circ h: B_3 \to B_2 \to A].
\end{align*}
So that for any map $g: B_1 \to A$, $F(h) \circ F(f)$ maps to $(g \circ h) \circ F(f) = g \circ h \circ f$, which is precisely $F(h \circ f)$, so that the functor is associative.
\end{proof} 

\begin{exercise}{12 (Micro)}
Show how 1.1.17 and 1.1.19 are special cases of 1.1.20
\end{exercise}
\begin{proof}
The way that 1.1.17 is a special case of 1.1.20 is more or less straightforward from the definition of the contravariant functor, since $(f \circ g)^\lor = g^\lor \circ f^\lor$ (here $g$ plays the same role as in example 1.1.20).
In this case, $A$ is $k$, the field in which the vector space is defined.

1.1.19 is a special case of 1.1.20 in a similar fashion.
We first note that $A$ in this case is $X$, and that if $f, g \in Hom(X, Y)$, then the contravariant functor maps $F(f \circ g)$ to $g \circ f$...
\end{proof} 

\begin{exercise}{13 (Micro)}
If $G$ and $H$ are groups, what functions $f : G \to H$ induce a functor between the one-object categories of $G$ and $H$.
\end{exercise}
\begin{proof}
If we define the categories as in exercise 9 (Micro), then we have that, if $f, g \in Hom(\ast, \ast) = G$, then $F(f \circ g) = g \circ f$;
that is, we invert the order of the operation.
\end{proof} 

\begin{exercise}{14 (Micro)}
If $P$ and $P'$ are posets, what functions $f: P \to P'$ induce a functor between the categories induced by $P$ and $P'$.
\end{exercise}
\begin{proof}
Using the formalism of exercise 10 (Micro), we have that for $x, y \in P$, we have that $F(x) \leq F(y)$ (that is, $Hom(F(y), F(x))$ is not empty) if and only if $x \geq y$.
In words, we invert the partial order.
These are the set of monotonically decreasing functions.
\end{proof} 
